% Options for packages loaded elsewhere
% Options for packages loaded elsewhere
\PassOptionsToPackage{unicode}{hyperref}
\PassOptionsToPackage{hyphens}{url}
\PassOptionsToPackage{dvipsnames,svgnames,x11names}{xcolor}
%
\documentclass[
  11pt,
  a4paper,
]{article}
\usepackage{xcolor}
\usepackage[top=2.5cm,bottom=2.5cm,left=2.5cm,right=2.5cm]{geometry}
\usepackage{amsmath,amssymb}
\setcounter{secnumdepth}{2}
\usepackage{iftex}
\ifPDFTeX
  \usepackage[T1]{fontenc}
  \usepackage[utf8]{inputenc}
  \usepackage{textcomp} % provide euro and other symbols
\else % if luatex or xetex
  \usepackage{unicode-math} % this also loads fontspec
  \defaultfontfeatures{Scale=MatchLowercase}
  \defaultfontfeatures[\rmfamily]{Ligatures=TeX,Scale=1}
\fi
\usepackage{lmodern}
\ifPDFTeX\else
  % xetex/luatex font selection
\fi
% Use upquote if available, for straight quotes in verbatim environments
\IfFileExists{upquote.sty}{\usepackage{upquote}}{}
\IfFileExists{microtype.sty}{% use microtype if available
  \usepackage[]{microtype}
  \UseMicrotypeSet[protrusion]{basicmath} % disable protrusion for tt fonts
}{}
\usepackage{setspace}
\makeatletter
\@ifundefined{KOMAClassName}{% if non-KOMA class
  \IfFileExists{parskip.sty}{%
    \usepackage{parskip}
  }{% else
    \setlength{\parindent}{0pt}
    \setlength{\parskip}{6pt plus 2pt minus 1pt}}
}{% if KOMA class
  \KOMAoptions{parskip=half}}
\makeatother


\usepackage{longtable,booktabs,array}
\usepackage{calc} % for calculating minipage widths
% Correct order of tables after \paragraph or \subparagraph
\usepackage{etoolbox}
\makeatletter
\patchcmd\longtable{\par}{\if@noskipsec\mbox{}\fi\par}{}{}
\makeatother
% Allow footnotes in longtable head/foot
\IfFileExists{footnotehyper.sty}{\usepackage{footnotehyper}}{\usepackage{footnote}}
\makesavenoteenv{longtable}
\usepackage{graphicx}
\makeatletter
\newsavebox\pandoc@box
\newcommand*\pandocbounded[1]{% scales image to fit in text height/width
  \sbox\pandoc@box{#1}%
  \Gscale@div\@tempa{\textheight}{\dimexpr\ht\pandoc@box+\dp\pandoc@box\relax}%
  \Gscale@div\@tempb{\linewidth}{\wd\pandoc@box}%
  \ifdim\@tempb\p@<\@tempa\p@\let\@tempa\@tempb\fi% select the smaller of both
  \ifdim\@tempa\p@<\p@\scalebox{\@tempa}{\usebox\pandoc@box}%
  \else\usebox{\pandoc@box}%
  \fi%
}
% Set default figure placement to htbp
\def\fps@figure{htbp}
\makeatother





\setlength{\emergencystretch}{3em} % prevent overfull lines

\providecommand{\tightlist}{%
  \setlength{\itemsep}{0pt}\setlength{\parskip}{0pt}}



 
\usepackage[style=authoryear-comp,]{biblatex}
\addbibresource{refs.bib}


\newcommand{\dt}{\text{d}}
\newcommand{\E}{\mathbb{E}}
\usepackage{mathtools}
\mathtoolsset{showonlyrefs,mathic=true}
\usepackage{orcidlink}
\definecolor{mypink}{RGB}{219, 48, 122}
\makeatletter
\@ifpackageloaded{tcolorbox}{}{\usepackage[skins,breakable]{tcolorbox}}
\@ifpackageloaded{fontawesome5}{}{\usepackage{fontawesome5}}
\definecolor{quarto-callout-color}{HTML}{909090}
\definecolor{quarto-callout-note-color}{HTML}{0758E5}
\definecolor{quarto-callout-important-color}{HTML}{CC1914}
\definecolor{quarto-callout-warning-color}{HTML}{EB9113}
\definecolor{quarto-callout-tip-color}{HTML}{00A047}
\definecolor{quarto-callout-caution-color}{HTML}{FC5300}
\definecolor{quarto-callout-color-frame}{HTML}{acacac}
\definecolor{quarto-callout-note-color-frame}{HTML}{4582ec}
\definecolor{quarto-callout-important-color-frame}{HTML}{d9534f}
\definecolor{quarto-callout-warning-color-frame}{HTML}{f0ad4e}
\definecolor{quarto-callout-tip-color-frame}{HTML}{02b875}
\definecolor{quarto-callout-caution-color-frame}{HTML}{fd7e14}
\makeatother
\makeatletter
\@ifpackageloaded{caption}{}{\usepackage{caption}}
\AtBeginDocument{%
\ifdefined\contentsname
  \renewcommand*\contentsname{Table of contents}
\else
  \newcommand\contentsname{Table of contents}
\fi
\ifdefined\listfigurename
  \renewcommand*\listfigurename{List of Figures}
\else
  \newcommand\listfigurename{List of Figures}
\fi
\ifdefined\listtablename
  \renewcommand*\listtablename{List of Tables}
\else
  \newcommand\listtablename{List of Tables}
\fi
\ifdefined\figurename
  \renewcommand*\figurename{Figure}
\else
  \newcommand\figurename{Figure}
\fi
\ifdefined\tablename
  \renewcommand*\tablename{Table}
\else
  \newcommand\tablename{Table}
\fi
}
\@ifpackageloaded{float}{}{\usepackage{float}}
\floatstyle{ruled}
\@ifundefined{c@chapter}{\newfloat{codelisting}{h}{lop}}{\newfloat{codelisting}{h}{lop}[chapter]}
\floatname{codelisting}{Listing}
\newcommand*\listoflistings{\listof{codelisting}{List of Listings}}
\usepackage{amsthm}
\theoremstyle{plain}
\newtheorem{theorem}{Theorem}[section]
\theoremstyle{definition}
\newtheorem{definition}{Assumption}[section]
\theoremstyle{plain}
\newtheorem{corollary}{Corollary}[section]
\theoremstyle{plain}
\newtheorem{lemma}{Lemma}[section]
\theoremstyle{remark}
\AtBeginDocument{\renewcommand*{\proofname}{Proof}}
\newtheorem*{remark}{Remark}
\newtheorem*{solution}{Solution}
\newtheorem{refremark}{Remark}[section]
\newtheorem{refsolution}{Solution}[section]
\makeatother
\makeatletter
\makeatother
\makeatletter
\@ifpackageloaded{caption}{}{\usepackage{caption}}
\@ifpackageloaded{subcaption}{}{\usepackage{subcaption}}
\makeatother
\usepackage{bookmark}
\IfFileExists{xurl.sty}{\usepackage{xurl}}{} % add URL line breaks if available
\urlstyle{same}
\hypersetup{
  pdftitle={Anomaly detection using surprisals},
  pdfauthor={Rob J Hyndman; David T Frazier},
  pdfkeywords={Extreme value theory, Generalized Pareto
Distribution, Outlier detection, Tail bounds},
  colorlinks=true,
  linkcolor={blue},
  filecolor={Maroon},
  citecolor={Blue},
  urlcolor={Blue},
  pdfcreator={LaTeX via pandoc}}

%% CAPTIONS
\usepackage{caption}
\DeclareCaptionStyle{italic}[justification=centering]
 {labelfont={bf},textfont={it},labelsep=colon}
\captionsetup[figure]{style=italic,format=hang,singlelinecheck=true}
\captionsetup[table]{style=italic,format=hang,singlelinecheck=true}

%% FONT
\usepackage{bera}
\usepackage[charter]{mathdesign}
%\usepackage[scale=0.9]{sourcecodepro}
\usepackage[lf,t]{FiraSans}

%% HEADERS AND FOOTERS
\usepackage{fancyhdr}
\pagestyle{fancy}
\rfoot{\Large\sffamily\raisebox{-0.1cm}{\textbf{\thepage}}}
\makeatletter
\lhead{\textsf{\expandafter{\@title}}}
\makeatother
\rhead{}
\cfoot{}
\setlength{\headheight}{15pt}
\renewcommand{\headrulewidth}{0.4pt}
\renewcommand{\footrulewidth}{0.4pt}
\fancypagestyle{plain}{%
\fancyhf{} % clear all header and footer fields
\fancyfoot[C]{\sffamily\thepage} % except the center
\renewcommand{\headrulewidth}{0pt}
\renewcommand{\footrulewidth}{0pt}}

%% MATHS
\usepackage{bm,amsmath}
\allowdisplaybreaks

%% GRAPHICS
\makeatletter
\def\fps@figure{htbp}
\makeatother
\setcounter{topnumber}{2}
\setcounter{bottomnumber}{2}
\setcounter{totalnumber}{4}
\renewcommand{\topfraction}{0.85}
\renewcommand{\bottomfraction}{0.85}
\renewcommand{\textfraction}{0.15}
\renewcommand{\floatpagefraction}{0.8}

%% SECTION TITLES
\usepackage[compact,sf,bf]{titlesec}
\titleformat{\section}[block]
  {\fontsize{15}{17}\bfseries\sffamily}
  {\thesection}
  {0.4em}{}
\titleformat{\subsection}[block]
  {\fontsize{12}{14}\bfseries\sffamily}
  {\thesubsection}
  {0.4em}{}
\titlespacing{\section}{0pt}{*5}{*1}
\titlespacing{\subsection}{0pt}{*2}{*0.2}
\titlespacing{\subsubsection}{0pt}{*1}{*0.1}

%% BIBLIOGRAPHY.

\makeatletter
\@ifpackageloaded{biblatex}{
\ExecuteBibliographyOptions{bibencoding=utf8,minnames=1,maxnames=3, maxbibnames=99,dashed=false,terseinits=true,giveninits=true,uniquename=false,uniquelist=false,doi=false, isbn=false,url=true,sortcites=false}
\DeclareFieldFormat{url}{\texttt{\url{#1}}}
\DeclareFieldFormat[article]{pages}{#1}
\DeclareFieldFormat[inproceedings]{pages}{\lowercase{pp.}#1}
\DeclareFieldFormat[incollection]{pages}{\lowercase{pp.}#1}
\DeclareFieldFormat[article]{volume}{\mkbibbold{#1}}
\DeclareFieldFormat[article]{number}{\mkbibparens{#1}}
\DeclareFieldFormat[article]{title}{\MakeCapital{#1}}
\DeclareFieldFormat[article]{url}{}
\DeclareFieldFormat[inproceedings]{title}{#1}
\DeclareFieldFormat{shorthandwidth}{#1}
\usepackage{xpatch}
\xpatchbibmacro{volume+number+eid}{\setunit*{\adddot}}{}{}{}
% Remove In: for an article.
\renewbibmacro{in:}{%
  \ifentrytype{article}{}{%
  \printtext{\bibstring{in}\intitlepunct}}}
\AtEveryBibitem{\clearfield{month}}
\AtEveryCitekey{\clearfield{month}}
\DeclareDelimFormat[cbx@textcite]{nameyeardelim}{\addspace}
\renewcommand*{\finalnamedelim}{\addspace\&\space}
}{}
\makeatother

%% PAGE BREAKING to avoid widows and orphans
\clubpenalty = 2000
\widowpenalty = 2000
\usepackage{microtype}
% Placement of logos

\RequirePackage[absolute,overlay]{textpos}
\setlength{\TPHorizModule}{1cm}
\setlength{\TPVertModule}{1cm}
\def\placefig#1#2#3#4{\begin{textblock}{.1}(#1,#2)\rlap{\includegraphics[#3]{#4}}\end{textblock}}

% Title and date

\title{Anomaly detection using surprisals}
\date{29 September 2025}

\def\Date{\number\day}
\def\Month{\ifcase\month\or
 January\or February\or March\or April\or May\or June\or
 July\or August\or September\or October\or November\or December\fi}
\def\Year{\number\year}

% Working paper number and JEL codes

\makeatletter
\def\wp#1{\gdef\@wp{#1}}\def\@wp{??/??}
\def\jel#1{\gdef\@jel{#1}}\def\@jel{??}
\def\showjel{{\large\textsf{\textbf{JEL classification:}}~\@jel}}
\def\nojel{\def\showjel{}}
\makeatother

\wp{no/yr}
\jel{C10,C14,C46}

% Title page

\makeatletter
\def\cover{{\sffamily\setcounter{page}{0}
        \thispagestyle{empty}
        \placefig{2}{1.5}{width=5cm}{monash2}
        \placefig{16.9}{1.5}{width=2.1cm}{MBSportrait}
        \begin{textblock}{4}(16.9,4)ISSN 1440-771X\end{textblock}
        \begin{textblock}{7}(12.7,27.9)\hfill
        \includegraphics[height=0.7cm]{AACSB}~~~
        \includegraphics[height=0.7cm]{EQUIS}~~~
        \includegraphics[height=0.7cm]{AMBA}
        \end{textblock}
        \vspace*{2.5cm}
        \begin{center}\Large
        Department of Econometrics and Business Statistics\\[.5cm]
        \footnotesize http://monash.edu/business/ebs/research/publications
        \end{center}\vspace{2cm}
        \begin{center}
        \fbox{\parbox{14cm}{\begin{onehalfspace}\centering\Huge\vspace*{0.3cm}
                \textsf{\textbf{\expandafter{\@title}}}\vspace{1cm}\par
                \LARGE
                \expandafter{\@author}
                \end{onehalfspace}
        }}
        \end{center}
        \vfill
                \begin{center}\Large
                \Month~\Year\\[1cm]
                Working Paper \@wp
        \end{center}\vspace*{2cm}}}
        \def\addresses#1{\gdef\@addresses{#1}}\def\@addresses{??}
        \def\pageone{{\sffamily\setstretch{1}%
        \thispagestyle{empty}%
        \vbox to \textheight{%
        \raggedright\baselineskip=1.2cm
     {\fontsize{24.88}{30}\sffamily\textbf{\expandafter{\@title}}}
        \vspace{2cm}\par
        \hspace{1cm}\parbox{14cm}{\sffamily\large\@addresses}\vspace{1cm}\vfill
        \hspace{1cm}{\large\Date~\Month~\Year}\\[1cm]
        \hspace{1cm}\showjel\vss}}}
\def\blindtitle{{\sffamily
     \thispagestyle{plain}\raggedright\baselineskip=1.2cm
     {\fontsize{24.88}{30}\sffamily\textbf{\expandafter{\@title}}}\vspace{1cm}\par
        }}
\def\titlepage{{\cover\newpage\pageone\newpage\blindtitle}}

\def\blind{\def\titlepage{{\blindtitle}}\let\maketitle\blindtitle}
\def\titlepageonly{\def\titlepage{{\pageone\end{document}}}}
\def\nocover{\def\titlepage{{\pageone\newpage\blindtitle}}\let\maketitle\titlepage}
\let\maketitle\titlepage
\makeatother

% Authors


  \author{Rob J Hyndman, David T Frazier}
  \addresses{%
    %
      \textbf{Rob J Hyndman}\\%
      %
        Department of Econometrics \& Business Statistics\\%
        Monash University\\%
        Clayton 3800\\%
        Australia\\%
      %
      {Email: Rob.Hyndman@monash.edu}\\%
      \textit{Corresponding author}\\[0.5cm]%
   %
      \textbf{David T Frazier}\\%
      %
        Department of Econometrics \& Business Statistics\\%
        Monash University\\%
        Clayton 3800\\%
        Australia\\%
      %
      %
      \\[0.5cm]%
   %
   }%
   \lfoot{\sf Hyndman, Frazier: 29 September 2025}

% Keywords

\newenvironment{keywords}{\par\vspace{0.5cm}\noindent{\sffamily\textbf{Keywords:}}}{\vspace{0.25cm}\par\hrule\vspace{0.5cm}\par}

% Abstract
\renewenvironment{abstract}{\begin{minipage}{\textwidth}\parskip=1.4ex\noindent
\hrule\vspace{0.1cm}\par{\sffamily\textbf{\abstractname}}\newline\setstretch{1.5}}
  {\end{minipage}}
\begin{document}
\maketitle

\begin{abstract}
TBC
\end{abstract}

\begin{keywords}
  Extreme value theory; Generalized Pareto Distribution; Outlier
detection; 
  Tail bounds.
\end{keywords}


\setstretch{1.5}
\section{Introduction}\label{introduction}

Anomaly detection methods are often rather ad hoc and lack a solid
theoretical foundation, or they make strong distributional assumptions
that may not hold in practice. Here we propose a general framework for
anomaly detection using extreme value theory to estimate the probability
of each observation. This approach allows for a unified treatment of
anomaly detection in both univariate and multivariate settings, and can
be applied to both unconditional and conditional distributions. The
basic idea is to calculate the surprisal (equal to the negative log
density) of each observation, and then to use extreme value theory to
model the tail of the surprisal distribution. This allows us to estimate
the probability of observing a surprisal value as extreme or more
extreme than a given observation, which can be used to flag anomalies.
The surprisal approach to anomaly detection provides a principled way to
control the false positive rate when flagging anomalies, by setting a
threshold on the estimated tail probabilities.

We show that under weak conditions on the tail of the surprisal
distribution, the maximum surprisal converges to one of the three types
of extreme value distributions. This justifies fitting a Generalized
Pareto Distribution (GPD) to the largest surprisal values, and using it
to estimate tail probabilities. We illustrate the approach with several
examples, including both simulated and real data.

Let \(\bm{y}_1,\dots,\bm{y}_n\) be iid observations on \(\mathbb{R}^d\)
from a probability distribution \(F\). The \emph{surprisal} of
observation \(\bm{y}_i\) is defined as \autocite{Stone2022} \[
s_i=-\log f(\bm{y}_i),
\] where \(f\) is either the density or probability mass function of
\(F\). While these are better known as ``log scores'' in statistics
\autocite{Good1952-zw}, in this context we prefer the term ``surprisal''
\autocite[coined by][]{Tribus61}, as it captures the idea that
observations with low density are more ``surprising'' (and therefore
potentially anomalous) than those with high density.

Let \(G(s) = \Pr\{S\le s\}\) be the distribution function of the
surprisal values, where \(S = -\log f(\bm{Y})\) and \(\bm{Y}\sim F\).
Then, the probability of observing a value with surprisal at least as
large as \(s\) is given by \(1-G(s)\), which leads to a natural way to
assign an anomaly score to each observation \(\bm{y}_i\): \[
p_i = 1-G(s_i) = 1-G(-\log f(\bm{y}_i)).
\] Hence, one way to detect anomalies is to flag observations with
\(p_i < \alpha\), giving a false positive rate of \(\alpha\) if all
observations are from \(F\). This is preferable to directly computing
the probabilities from the tail of \(F\), as it allows for anomalies to
occur anywhere in the distribution of \(F\), not just in its tails. For
example, anomalies may lie between modes in a multimodal distribution,
or in otherwise low-density regions of the distribution \(F\). It is
also preferable to a distance-based approach to anomaly detection
\autocite[e.g.,][]{hdoutliers}, as it naturally handles skewed and
heavy-tailed distributions.

The surprisal approach converts the problem of working with a
potentially multivariate distribution \(F\) to working with a univariate
distribution \(G\). We can write \[
G(s) = \Pr\{-\log f(Y)\le s\} = \Pr\{ f(Y)\ge e^{-s}\}.
\] This makes clear that the surprisal probability is related to the
highest density regions of the distribution \autocite{HDR96}. In fact,
every observation with \(p_i< \alpha\) lies outside the
\(100(1-\alpha)\)\% HDR region of \(F\).

We will usually not know \(F\), and so we need to estimate it from data
or from theory. Such estimates are particularly fragile in areas of low
density, which is precisely where we are looking for anomalies.

To address this, \textcite{lookout} proposed estimating \(F\) by a
leave-one-out kernel density estimate, and then fitting a Generalized
Pareto Distribution (GPD) to the most extreme surprisal values. The
underlying idea is that even if \(F\) is not particularly
well-estimated, the tail of the surprisal distribution should still be
well-approximated by a GPD. \textcite{lookout2} provides some
theoretical justification for this approach in the context of kernel
density estimation.

Here we extend that idea to allow \(F\) to be any well-behaved
distribution, including conditional distributions. Specifically, we fit
a GPD to the largest \(\beta\) of \(\{s_i\}_{i=1}^n\), obtaining
estimates of the location, scale and shape parameters,
\((\mu,\sigma,\xi)\). We calculate the probability of each observation
or one more extreme,
\(p_i = \beta\left[1-P(s_i \mid \hat\mu, \hat\sigma, \hat\xi)\right]\),
where \(P\) is the distribution function of the GPD. Observations not
contained in the largest \(\beta\) of \(\{s_i\}_{i=1}^n\) are assigned
\(p_i = \beta\).

To justify this approach, we seek conditions on the surprisal
distribution of \(S = -\log f(\bm{Y})\), where \(\bm{Y}\sim F\), similar
to the Fisher-Tippett-Gnedenko (FTG) theorem
\autocite{coles2001introduction}.

\begin{tcolorbox}[enhanced jigsaw, bottomtitle=1mm, colbacktitle=quarto-callout-note-color!10!white, opacitybacktitle=0.6, bottomrule=.15mm, colframe=quarto-callout-note-color-frame, opacityback=0, coltitle=black, title=\textcolor{quarto-callout-note-color}{\faInfo}\hspace{0.5em}{Note}, colback=white, toptitle=1mm, titlerule=0mm, leftrule=.75mm, arc=.35mm, breakable, left=2mm, rightrule=.15mm, toprule=.15mm]

Outline of paper TBC

\end{tcolorbox}

\section{Extreme value theory for the surprisal
distribution}\label{extreme-value-theory-for-the-surprisal-distribution}

We consider three possible conditions on the surprisal distribution of
\(S = -\log f(\bm{Y})\) where \(f\) is the density or mass function of
\(\bm{Y}\). These conditions describe increasingly heavy tails:

\begin{enumerate}
\def\labelenumi{\arabic{enumi}.}
\tightlist
\item
  where \(S\) is sub-Gaussian (e.g., when \(f\) is bounded above and
  below);
\item
  where \(S\) is sub-exponential (e.g., when \(f\) is Gaussian); and
\item
  where \(S\) has a polynomial moment (e.g., when \(f\) is Student-t).
\end{enumerate}

We state these as three alternative assumptions.

\begin{definition}[Sub-Gaussian]\protect\hypertarget{def-bounded}{}\label{def-bounded}

The random variable \(S =-\log f(Y)\in\mathbb{R}\), and satisfies, for
all \(\lambda\in\mathbb{R}\), and some \(\nu>0\), \[
    \E \exp\{\lambda (S-\E[S])\}\le \exp\{\lambda^2 \nu^2/2\}.
\]

\end{definition}

\begin{definition}[Sub-exponential]\protect\hypertarget{def-subbounded}{}\label{def-subbounded}

The random variable \(S\) is sub-exponential with parameters \(\nu\) and
\(b\), i.e., \[
    \E \exp\{\lambda(S-\E[S])\} \leq \exp\{\lambda^2\nu^2/ 2\} \quad\text{ for all } \quad |\lambda|< 1/b.
\]

\end{definition}

\begin{definition}[Polynomial]\protect\hypertarget{def-polybounded}{}\label{def-polybounded}

The random variable \(|S|\) has polynomial moments of order
\(\alpha\ge 1\); i.e., \(\mathbb{E}[|S|^p] \le C^p\) for some \(C>0\)
such that \(C^p-1>0\).

\end{definition}

Under each of Assumptions \ref{def-bounded}--\ref{def-polybounded}, we
seek a result on the limiting distribution of the maximum of the
surprisals \(M_n = \max\{s_1,\dots,s_n\}\). If we can show that this
distribution falls within the class of extreme value distributions, then
we can fit a Generalized Pareto Distribution (GPD) to the largest
surprisal values \autocite{coles2001introduction}.

\subsection{\texorpdfstring{\(S\) has sub-Gaussian
tails}{S has sub-Gaussian tails}}\label{s-has-sub-gaussian-tails}

Let's first assume that the surprisals are well-behaved and have a
sub-Gaussian tail. Assumption~\ref{def-bounded} will be satisfied in
cases where the density \(f\) is bounded both above and below, such as
with mass functions on a finite set. However, it may not be satisfied if
the log density has no lower bound.

For example, consider the truncated normal distribution on \([a,b]\),
where \(Y\sim \mathcal{TN}(a,b,\mu,\sigma^2)\) and
\(-\infty < a < b < \infty\). Then \[
S=-\log(2\pi\sigma^2)/2-\frac{1}{2\sigma^2}(Y-\mu)^2-\log\{\Phi(b)-\Phi(a)\},
\] where \(\Phi\) is the standard normal distribution function. Since
\(Y\in[a,b]\), \(S\) is bounded for all \(\sigma>0\) and hence \(S\)
satisfies Assumption~\ref{def-bounded}.

Under Assumption~\ref{def-bounded}, we have the following result on the
maximum of the surprisals.

\begin{lemma}[]\protect\hypertarget{lem-bounds}{}\label{lem-bounds}

Suppose \(\bm{y}_1,\dots,\bm{y}_n\) are iid from a density \(f\), and
let \(M_n = \max\{s_1,\dots,s_n\}\) where \(s_i = -\log f(\bm{y}_i)\).
If \(S =-\log f(Y)\in\mathbb{R}\), and satisfies, for all
\(\lambda\in\mathbb{R}\), and some \(\nu>0\), \[
\mathbb{E} \exp\{\lambda (S-\mathbb{E}[S])\}\le \exp\{\lambda^2 \nu^2/2\},
\], then, for \(s\in\mathbb{R}_+\), \[
\Pr\{M_n-\mathbb{E}[S]\ge s\}\le 1-\left(1-e^{-\frac{1}{2}s^2/\nu^2}\right)^n.
\]

\end{lemma}

\begin{proof}
For any \(\lambda>0\), \begin{align}
\Pr\{S\ge s\} =
\Pr\left\{e^{\lambda (S-\mathbb{E}[S])}\ge e^{\lambda (s-\mathbb{E}[S])}\right\}
 & \le e^{-\lambda(s-\mu)}\mathbb{E} e^{\lambda (S-\mu)} \\
 & \le e^{-\lambda(s-\mu)}e^{\lambda^2\nu^2/2} .
\end{align} Now, we can maximise the RHS of the above bound in
\(\lambda\), which yields the optimal value \(\lambda=(s-\mu)/\{M^2\}\).
Plugging this in we have the bound \[
\Pr\{S\ge s\}\le e^{-\frac{1}{2}(s-\mu)^2/\nu^2}.
\] The stated result now follows by noting that \begin{align}
\Pr\left(M_n-\mathbb{E}[S]\ge s\right) = 1-\Pr\{M_n\le s\}
 & = 1-\Pr\{S\le s\}^n                       \\
 & = 1-(1-\Pr\{S\ge s\})^n                           \\
 & \le 1-\left(1-e^{-\frac{1}{2}s^2/\nu^2}\right)^n.
\end{align}
\end{proof}

\begin{corollary}[]\protect\hypertarget{cor-one}{}\label{cor-one}

If Assumption~\ref{def-bounded} is satisfied, then, for all \(s>0\), \[
\Pr\Big(|M_n-\mathbb{E}[S]| \ge \sqrt{2\nu^2s} + \sqrt{2\nu^2\log(2n)}\Big) \le e^{-s}.
\]

\end{corollary}

\begin{proof}
This follows directly from the union bound: \begin{align}
\Pr\Big(|M_n|\ge \sqrt{2\nu^2s} + \sqrt{2\nu^2\log(n)}\Big)
 & \le n \Pr\Big\{|S|\ge \sqrt{2}\nu^2s + \sqrt{2\nu^2\log(n)}\Big\} \\
 & \le 2n e^{-[\sqrt{2\nu^2s}+\sqrt{2\nu^2\log(n)}]^2/(2\nu^2)} \\
 & = 2n e^{-\frac{1}{2\nu^2}(2s\nu^2 )-\frac{2\nu^2\log(2n)}{2\nu^2}-\frac{\sqrt{2\nu^2s}\sqrt{2\nu^2\log(2n)}}{2\nu^2}} \\
 & \le e^{-\sqrt{\log(2n)s}-s} \\
 & \le e^{-s}
\end{align} The first result follows from the iid assumption on \(S\)
and the union bound, while the second from the assumption that \(S\) is
sub-Gaussian. The remaining inequalities follow from reorganising terms.
\end{proof}

Corollary~\ref{cor-one} is useful as it clarifies that if the surprisal
is sub-Gaussian, then the maximal surprisal will follow a reversed
Weibull distribution, and thus is within the class of extreme value
distributions.

\subsection{\texorpdfstring{\(S\) has an exponential
moment}{S has an exponential moment}}\label{s-has-an-exponential-moment}

In cases where the underlying density has full support on \(\mathbb{R}\)
and the log density cannot be bounded from below, then the
sub-Gaussianity condition in Assumption~\ref{def-bounded} may be
violated, as the following result illustrates.

\begin{lemma}[]\protect\hypertarget{lem-normal}{}\label{lem-normal}

If \(Y\sim \mathcal{N}(\mu,\sigma^2)\), then \(S=-\log f(Y)\) is
sub-exponential. In particular, \[
\mathbb{E} e^{\lambda(S-1/2)+(\lambda/2)\log(2\pi\sigma^2)}\le e^{\lambda^2/2+(\lambda/2)\log(2\pi\sigma^2)},\quad \forall |\lambda|<1/2.
\]

\end{lemma}

\begin{proof}
Firstly, note that \[
S=-\log f(Y)= \frac{1}{2}\log (2\pi\sigma^2)+\frac{(Y-\mu)^2}{2\sigma^2}.
\] and that \[
\mu_S=\mathbb{E}[S]=-\int f(y)\log f(y)\mathrm{d} y=\frac{1}{2}\{1+\log (2\pi\sigma^2)\}.
\] Define \[
X=\frac{1}{2}\frac{(Y-\mu)^2}{\sigma^2}-\frac{1}{2},
\] then \(\mathbb{E}[X]=0\). Now, note that \[
\mathbb{E} e^{\frac{\lambda}{2}(S-1)}={e^{\lambda/2\{\log(2\pi\sigma^2)\}}}\frac{1}{\sqrt{2\pi\sigma^2}}\int e^{\frac{\lambda}{2}\left[\frac{(y-\mu)^2}{\sigma^2}-1\right]}e^{-\frac{(y-\mu)^2}{2\sigma^2}}\mathrm{d} y.
\] Consider the change of variables \(z=(y-\mu)/\sigma\) so that we have
\[
\E e^{\frac{\lambda}{2}(S-1)}={e^{\lambda/2\{\log(2\pi\sigma^2)\}}}\frac{1}{\sqrt{2\pi}}\int e^{\frac{\lambda}{2}[z^2-1]}e^{-z^2/2}\dt y={e^{\lambda/2\{\log(2\pi\sigma^2)\}}}\frac{e^{-2(\lambda/2)}}{\sqrt{1-2(\lambda/2)}}
\] for all \(\lambda<1\). Further, it can be verified that \[
\frac{e^{-\lambda/2}}{\sqrt{1-2(\lambda/2)}}\le e^{\lambda^2/2},\quad \forall |\lambda|<1/2.
\] Hence, when we have \(Y\sim\mathcal{N}(\mu,\sigma^2)\), then
\(S=-\log f(Y)\) is sub-exponential with parameter \(\nu=1\) and
\(b=2\).
\end{proof}

Therefore, to be able to apply extreme value theory to the surprisal
values, we need a version of Lemma~\ref{lem-bounds} that uses
sub-exponentiality.

\begin{lemma}[]\protect\hypertarget{lem-bounds_sub}{}\label{lem-bounds_sub}

Suppose \(\bm{y}_1,\dots,\bm{y}_n\) are iid from a density \(f\), and
let \(M_n = \max\{s_1,\dots,s_n\}\) where \(s_i = -\log f(\bm{y}_i)\).
If Assumption~\ref{def-subbounded} is satisfied, then, for
\(s\in\mathbb{R}_+\), \[
    \Pr\left(M_n-\E[S]\ge s\right)\le 1-
    \begin{cases}
        (1-e^{-\frac{1}{2}s^2/\nu^2})^n & \text{ if }0\le s\le \nu^2/b \\
        (1-e^{-s/2b})^n                 & \text{ if }s>\nu^2/b
    \end{cases}.
\]

\end{lemma}

\begin{proof}
The result follows a similar argument as for Lemma~\ref{lem-bounds}, up
until we maximize the tail bound. In particular, for any \(\lambda>0\),
\begin{align*}
  \Pr\left\{ S-\E[S]\ge s\right\} =
  \Pr\left\{e^{\lambda (S-\E[S])}\ge e^{\lambda s}\right\}
    & \le e^{-\lambda s}\E e^{\lambda (S-\E[S])} \\
    & \le e^{-\lambda s}e^{\lambda^2\nu^2/2} .
\end{align*} Now, we must consider a constrained optimization to ensure
the right hand side exists. In particular, we know from the proof of
Lemma~\ref{lem-bounds} that the unconstrained optimum is
\(\lambda=s/\nu^2\), which corresponds to the constrained optimum on the
set \(0\le t<\nu^2/b\). However, when \(s\ge \nu^2/b\), the maximum is
achieved at \(\lambda=1/b\), and putting this together yields \[
  \Pr\left\{S-\E[S]\ge s\right\}\le
  \begin{cases}
    e^{-\frac{1}{2}s^2/\nu^2} & \text{ if }0\le s\le \nu^2/b; \\
    e^{-s/2b}                 & \text{ if } s>\nu^2/b.
  \end{cases}
\] The stated result again follows by noting that \begin{align*}
  \Pr\left\{M_n-\E[S]\ge s\right\}
    & = 1-\Pr\left\{M_n\le s\right\} \\
    & = 1-\Pr\left\{S\le s\right\}^n \\
    & \le 1-
  \begin{cases}
    (1-e^{-\frac{1}{2}s^2/\nu^2})^n & \text{ if } 0\le s\le \nu^2/b; \\
    (1-e^{-s/2b})^n                 & \text{ if } s>\nu^2/b.
  \end{cases}
\end{align*}
\end{proof}

Results similar to Lemma~\ref{lem-bounds_sub} also hold for
\(\Pr\left\{M_n-\E[S]\le s\right\}\), and for
\(\Pr\left\{|M_n-\E[S]|\ge s\right\}\) with an additional factor of
\(2\), i.e., \[
\Pr\left\{|M_n-\E[S]|\ge s\right\}\le 1-
\begin{cases}
  (1-2e^{-\frac{1}{2}s^2/\nu^2})^n & \text{ if }0\le s\le \nu^2/b \\
  (1-2e^{-s/2b})^n                 & \text{ if }s>\nu^2/b.
\end{cases}
\]

\begin{corollary}[]\protect\hypertarget{cor-two}{}\label{cor-two}

Under Assumption~\ref{def-subbounded}, for \(s>\nu^2/b\), as
\(n\rightarrow\infty\) \[
    \Pr\left(|M_n-\E[S]|\ge (2b)s+(2b)\log(2n)\right)\le 1-e^{-e^{-s}}+o(1/n).
\]

\end{corollary}

\begin{proof}
Apply Assumption~\ref{def-subbounded} to see that \begin{align*}
    \Pr\left(|M_n-\E[S]|\ge (2b)s+(2b)\log(2n)\right)
     & \le 1- \left\{1-2e^{-[2bs+2b\log(2n)]/2b}\right\}^n                       \\
     & = 1-\big(1-e^{-s}/n\big)^n                                     \\
     & = 1-e^{e^{-s}}-\left[\big(1-e^{-s}/n\big)^n-e^{e^{-s}}\right].
\end{align*} The result now follows by noting that, for all \(s>0\), \[
    \left[\big(1-e^{-s}/n\big)^n-e^{e^{-s}}\right]= o(1)\le 0.
\]
\end{proof}

Corollary~\ref{cor-two} is useful as it shows that if the surprisal has
a sub-exponential tail, then the resulting maximum surprisal tail
probability converges to that of an extreme value distribution. In
particular, we have that Corollary~\ref{cor-two} implies \[
    \sup_{s:s>\nu^2/b}\left|\Pr\Big\{|M_n-\mu|\ge \big[2bs+2b\log(2n)\big]\Big\}-\big(1-e^{e^{-s}}\big)\right| = o(1).
\] To see that this implies that \(M_n-\E[S]\), appropriately
standardized, converges to an extreme value distribution, we first
recall that \[
    \Pr\{X_n>x\}-\Pr\{X>x\}=1-\Pr\{X_n\le x\}-\left[1-\Pr\{X\le x\}\right]=\Pr\{X\le x\}-\Pr\{X_n\le x\}.
\] Hence, \[
    \sup_{x\in\mathcal{X}}|\Pr\{X_n\le x\}-\Pr\{X\le x\}|=\sup_{x\in\mathcal{X}}|\Pr\{X_n> x\}-\Pr\{X> x\}|,
\] so that when \(X_n\) and \(X\) are both continuous we see that
Corollary~\ref{cor-two} implies that the CDF of
\((M_n-\E[S]-2b\log(2n))/2b\) converges to that of the Type-1 extreme
value distribution (i.e., Gumbel distribution); i.e., using
Corollary~\ref{cor-two} and the above we have that \[
    \sup_{s:s>\nu^2/b}\left|\Pr\left(|M_n-\mu|\le [2bs+2b\log(2n)]\right\}-e^{e^{-s}}\right|=o(1).
\] This is one notion of convergence in distribution.

\subsection{\texorpdfstring{\(S\) has a polynomial
moment}{S has a polynomial moment}}\label{s-has-a-polynomial-moment}

Finally, we consider the case where the surprisal has a polynomial
moment (e.g., Student-t random variables), and derive a corresponding
result to Lemma~\ref{lem-bounds}.

Suppose \(\bm{y}_1,\dots,\bm{y}_n\) are iid from a density \(f\), and
let \(M_n = \max\{s_1,\dots,s_n\}\) where \(s_i = -\log f(\bm{y}_i)\).
If Assumption~\ref{def-polybounded} is satisfied, then, for \(s>0\) and
such that \(s>C>0\), \[
  \Pr\left(|M_n-\E[S]|\ge s\right) \le 1- (1-C^p/s^p)^n.
\] Hence, \[
    \sup_{s:s>C} \left|\Pr\big(|M_n-\E[S]|\le (s\cdot n^{1/p})\big)-e^{-s^{-p}}\right|=o(1).
\]

\begin{proof}
By Markov's inequality, \[
    \Pr\left\{ |S-\E[S]|\ge s\right\}\le \frac{\E|S|^p}{s^p}\le \frac{C^p}{s^p}.
\] To derive the first stated result, again note that \begin{align*}
\Pr\left\{|M_n-\E[S]|\ge s\right\}
& = 1 -\Pr\left\{|M_n-\E[S]|\le s\right\}               \\
& = 1-\Pr\left\{|S-\E[S]|\le s\right\}^n                \\
& = 1-\left(1-\Pr\left\{|S-\E[S]|\ge s\right\}\right)^n \\
& \le 1-(1-C^p/s^p)^n.
\end{align*}

To derive the second stated result, apply the first result with
\(t=sn^{1/p}C>1\), which clearly holds for some \(n\) large enough, to
see that \[
    \Pr\left\{|M_n-\E[S]|\ge t\right\}
    \le 1- \left(1-\frac{C^p}{(Csn^{1/p})^p}\right)^n \\
    = 1-\left(1-s^{-p}/n\right)^n.
\] Now, recall that, for \(x\ge0\), \[
    \lim_{n\rightarrow\infty}(1-x/n)^n=e^{-x}.
\] Since \(x\mapsto e^{-x}\) is bounded for all \(x\ge0\), it follows
that \[
    \sup_{s:s>C}\left|\Pr\left\{|M_n-\E[S]|\ge t\right\}-e^{-s^{-p}}\right|=o(1).
\]
\end{proof}

\subsection{Three types theorem}\label{three-types-theorem}

Placing the three lemmas together gives the following ``three types''
Theorem.

\begin{theorem}[]\protect\hypertarget{thm-three_types}{}\label{thm-three_types}

Let \(\bm{y}_1,\dots,\bm{y}_n\) be an iid sequence from \(F\), and let
\(S =-\log f(\bm{Y})\), where \(Y\sim \bm{F}\). Further, let
\(M_n=\max\{s_1,\dots,s_n\}\), where \(s_i=-\log f(y_i)\).

\begin{enumerate}
\def\labelenumi{\arabic{enumi}.}
\tightlist
\item
  If the surprisal satisfies Assumption~\ref{def-bounded}, then \[
    \sup_{s:s>0}\left|\Pr\left\{\Big|M_n-\E[S]\Big|\ge \sqrt{2\nu^2 s}+\sqrt{2\nu^2\log(2n)}\right\}-e^{-s}\right|=o(1).
  \]
\item
  If the surprisal satisfies Assumption~\ref{def-subbounded}, then \[
    \sup_{s:s>1/b}\left|\Pr\Big\{\big|M_n-\E[S]\big|\ge (2b)s+(2b)\log(2n)\Big\}-e^{-e^{-s}}\right|=o(1).
  \]
\item
  If the surprisal satisfies Assumption~\ref{def-polybounded}, then \[
   \sup_{s:s>C}\left|\Pr\Big\{\big|M_n-\E[S]\big|\ge (Csn^{1/p})\Big\}-e^{-s^{-p}}\right|=o(1).
  \]
\end{enumerate}

\end{theorem}

Theorem~\ref{thm-three_types} gives a corresponding `three types'
theorem for the behavior of surprisals and has a simple interpretation.
If the surprisal has Gaussian like tails, then the maximum surprisal is
a reversed Weibull; if it only has an exponential moment, then the
maximum surprisal will be Gumbel; if the maximum surprisal only has a
polynomial moment, then it is Fréchet. As each of these distributions is
a member of the generalized extreme value family, then the upper tail of
the distribution can be approximated by a Generalized Pareto
Distribution \autocite[Theorem 4.1]{coles2001introduction}.

\section{Experiments}\label{experiments}

\begin{itemize}
\tightlist
\item
  990 from Gamma(2,2) and 10 from Gamma(2,r). Fit Gamma to all 1000 and
  compute surprisals. Fit GPD to largest 50 surprisals. Compare true and
  estimated tail probabilities.
\item
  990 from trivariate N(0,1) with 10 from ring at distance 2 from
  origin. Fit bivariate normal to all 1000 and compute surprisals. Fit
  GPD to largest 50 surprisals. Compare true and estimated tail
  probabilities.
\item
  \url{https://otexts.com/weird/06-density.html\#example-standard-normal-values}
\end{itemize}

\section{Applications}\label{applications}

\begin{itemize}
\tightlist
\item
  US COVID death rates:
  \url{https://otexts.com/weird/09-timeseries.html\#multivariate-time-series}
\item
  \url{https://otexts.com/weird/06-density.html\#gam-surprisals}
\end{itemize}

\section{Discussion}\label{discussion}

\subsection*{Software and
reproducibility}\label{software-and-reproducibility}
\addcontentsline{toc}{subsection}{Software and reproducibility}

All results presented here can be reproduced using the code available at
\url{https://github.com/robjhyndman/Surprisal_Theory}. The analysis was
conducted using R version 4.5.1 \autocite{R}, with the following R
packages: weird \autocite{Rweird}, tsibble \autocite{Rtsibble}, targets
\autocite{targets,targetspaper}, ggplot2 \autocite{ggplot2,ggplotbook},
and other tidyverse \autocite{tidyverse} packages.

\subsection*{Acknowledgements}\label{acknowledgements}
\addcontentsline{toc}{subsection}{Acknowledgements}

Both authors are funded by the Australian Research Council through
Discovery Project DP250100702. Rob Hyndman is a member of the Australian
Research Council Industrial Transformation Training Centre in
Optimisation Technologies, Integrated Methodologies, and Applications
(OPTIMA), Project ID IC200100009.


\printbibliography



\end{document}
